\subsection{Definition and first properties of the torsion product}

For every A-module E, we denote by \(p_{E}: L(E) \to E\) the canonical free resolution of E (X, p.~50).

DEFINITION 1. --- Let M be a right A-module and N a left A-module. The torsion product of M and N is the graded k-module

\[
\operatorname{Tor} ^ {\mathbf {A}} (\mathbf {M}, \mathbf {N}) = \mathrm{H} (\mathrm{L} (\mathbf {M}) \otimes_ {\mathbf {A}} \mathrm{L} (\mathbf {N}))  .\tag{4}
\]

The homogeneous components of Tor \(^{A}\) (M, N) are denoted

\[
\operatorname{Tor} _ {n} ^ {\mathbf {A}} (\mathbf {M}, \mathbf {N}) = \mathrm{H} _ {n} (\mathrm{L} (\mathbf {M}) \otimes_ {\mathbf {A}} \mathrm{L} (\mathbf {N})).\tag{5}
\]

Since L(M) and L(N) are zero on the right, we have

\[
\operatorname{Tor} _ {n} ^ {\mathbf {A}} (\mathbf {M}, \mathbf {N}) = 0 \quad \text { for } \quad n <   0.\tag{6}
\]

Remark 1. --- We shall see below (X, p.~107, prop. 6) finiteness properties of the Tor \(^{A}\) (M, N) modules. For example, if A is commutative Noetherian and M and N are finitely generated A-modules, then each A-module Tor \(_{n}^{A}\) (M, N) is finitely generated.

Let \(f: \mathbf{M} \to \mathbf{M}'\) be a homomorphism of right A-modules and \(g: \mathbf{N} \to \mathbf{N}'\) be a homomorphism of left A-modules; we set \(\operatorname{Tor}^{\mathbf{A}}(f, g) = \operatorname{H}\left(\operatorname{L}(f) \otimes_{\mathbf{A}} \operatorname{L}(g)\right)\) ; it is a homomorphism of graded \(k\) -modules

\[
\operatorname{Tor} ^ {\mathbf {A}} (f, g): \operatorname{Tor} ^ {\mathbf {A}} (\mathbf {M}, \mathbf {N}) \rightarrow \operatorname{Tor} ^ {\mathbf {A}} (\mathbf {M} ^ {\prime}, \mathbf {N} ^ {\prime})
\]

whose homogeneous components are denoted

\[
\operatorname{Tor} _ {n} ^ {\mathbf {A}} (f, g): \operatorname{Tor} _ {n} ^ {\mathbf {A}} (\mathbf {M}, \mathbf {N}) \rightarrow \operatorname{Tor} _ {n} ^ {\mathbf {A}} (\mathbf {M} ^ {\prime}, \mathbf {N} ^ {\prime}).
\]

By prop. 1 of X, p.~62, the canonical homomorphism

\[
\gamma_ {0, 0}: \mathrm{H} _ {0} (\mathrm{L} (\mathrm{M})) \otimes_ {\mathrm{A}} \mathrm{H} _ {0} (\mathrm{L} (\mathrm{N})) \rightarrow \mathrm{H} _ {0} (\mathrm{L} (\mathrm{M}) \otimes_ {\mathrm{A}} \mathrm{L} (\mathrm{N}))
\]

is bijective; using the isomorphisms \(\mathbf{M} \to \mathbf{H}_{0}(\mathbf{L}(\mathbf{M}))\) and \(\mathbf{N} \to \mathbf{H}_{0}(\mathbf{L}(\mathbf{N}))\) we derive an isomorphism, called canonical

\[
\gamma_ {\mathbf {M}, \mathbf {N}}: \mathbf {M} \otimes_ {\mathbf {A}} \mathbf {N} \rightarrow \operatorname{Tor} _ {0} ^ {\mathbf {A}} (\mathbf {M}, \mathbf {N}).\tag{7}
\]

We shall always identify \(\operatorname{Tor}_{0}^{\mathbf{A}}(\mathbf{M},\mathbf{N})\) with \(\mathbf{M}\otimes_{\mathbf{A}}\mathbf{N}\) by this isomorphism. Then the \(k\) -linear map \(\operatorname{Tor}_{0}^{\mathbf{A}}(f,g)\) is identified with \(f\otimes g\) .

Remark 2. --- The morphism of complexes \(p_{\mathbf{M}} \otimes p_{\mathbf{N}}: \mathbf{L}(\mathbf{M}) \otimes_{\mathbf{A}} \mathbf{L}(\mathbf{N}) \to \mathbf{M} \otimes_{\mathbf{A}} \mathbf{N}\) induces on the homology of degree 0 the isomorphism

\[
\gamma_ {\mathbf {M}, \mathbf {N}} ^ {- 1}: \operatorname{Tor} _ {0} ^ {\mathbf {A}} (\mathbf {M}, \mathbf {N}) \rightarrow \mathbf {M} \otimes_ {\mathbf {A}} \mathbf {N}
\]

inverse of \(\gamma_{M,N}\) .

One has \(L(1_{\mathbf{M}}) = 1_{L(\mathbf{M})}\) , \(L(1_{\mathbf{N}}) = i_{L(\mathbf{N})}\) hence by passing to homology:

\[
\operatorname{Tor} ^ {\mathbf {A}} \left(1 _ {\mathbf {M}}, 1 _ {\mathbf {N}}\right) = 1 _ {\operatorname{Tor} ^ {\mathbf {A}} (\mathbf {M}, \mathbf {N})}.\tag{8}
\]

If \(f': \mathbf{M}' \to \mathbf{M}''\) (resp. \(g': \mathbf{N}' \to \mathbf{N}''\) ) is a homomorphism of right (resp. left) A-modules, we have \(\mathbf{L}(g' \circ g) = \mathbf{L}(g') \circ \mathbf{L}(g)\) and \(\mathbf{L}(f' \circ f) = \mathbf{L}(f') \circ \mathbf{L}(f)\) , hence

\[
\operatorname{Tor} ^ {\mathbf {A}} \left(f ^ {\prime} \circ f, g ^ {\prime} \circ g\right) = \operatorname{Tor} ^ {\mathbf {A}} \left(f ^ {\prime}, g ^ {\prime}\right) \circ \operatorname{Tor} ^ {\mathbf {A}} (f, g).\tag{9}
\]

Consider the morphisms of k-complexes

\[
\mathrm{L} (\mathrm{M}) \otimes_ {\mathrm{A}} \mathrm{N} \xleftarrow {1 \otimes p _ {\mathrm{N}}} \mathrm{L} (\mathrm{M}) \otimes_ {\mathrm{A}} \mathrm{L} (\mathrm{N}) \xrightarrow {p _ {\mathrm{M}} \otimes 1} \mathrm{M} \otimes_ {\mathrm{A}} \mathrm{L} (\mathrm{N})
\]

and the k-homomorphisms they induce in homology:

\[
\mathrm{H} (\mathrm{L} (\mathrm{M}) \otimes_ {\mathrm{A}} \mathrm{N}) \xleftarrow {\psi_ {\mathrm{M}} (\mathrm{N})} \operatorname{Tor} ^ {\mathrm{A}} (\mathrm{M}, \mathrm{N}) \xrightarrow {\overline {{\psi}} _ {\mathrm{N}} (\mathrm{M})} \mathrm{H} (\mathrm{M} \otimes_ {\mathrm{A}} \mathrm{L} (\mathrm{N}));
\]

by prop. 4 of X, p.~67, \(1 \otimes p_{\mathbf{N}}\) and \(p_{\mathbf{M}} \otimes 1\) are homologisms. Hence:

Proposition 5. --- The k-homomorphisms

\[
\begin{array}{l} \psi_ {\mathrm{M}} (\mathrm{N}): \operatorname{Tor} ^ {\mathrm{A}} (\mathrm{M}, \mathrm{N}) \to \mathrm{H} (\mathrm{L} (\mathrm{M}) \otimes_ {\mathrm{A}} \mathrm{N}) \\ \overline {{\psi}} _ {\mathrm{N}} (\mathrm{M}): \operatorname{Tor} ^ {\mathrm{A}} (\mathrm{M}, \mathrm{N}) \to \mathrm{H} (\mathrm{M} \otimes_ {\mathrm{A}} \mathrm{L} (\mathrm{N})) \end{array}
\]

are bijective.

COROLLARY. --- If M or N is flat, \(\operatorname{Tor}_{i}^{\mathbf{A}}(\mathbf{M},\mathbf{N}) = 0\) for \(i \geqslant 0\) .

Suppose N (resp. M) is flat; then \(p_{\mathbf{M}} \otimes 1 : \mathbf{L}(\mathbf{M}) \otimes_{\mathbf{A}} \mathbf{N} \to \mathbf{M} \otimes_{\mathbf{A}} \mathbf{N}\) (resp. \(1 \otimes p_{\mathbf{N}} : \mathbf{M} \otimes_{\mathbf{A}} \mathbf{L}(\mathbf{N}) \to \mathbf{M} \otimes_{\mathbf{A}} \mathbf{N}\) ) is a homologism (X, p.~67, prop. 4), hence \(\mathrm{H}_{i}(\mathrm{L}(\mathrm{M}) \otimes_{\mathbf{A}} \mathrm{N})\) (resp. \(\mathrm{H}_{i}(\mathrm{M} \otimes_{\mathbf{A}} \mathrm{L}(\mathrm{N}))\) ) is zero for i \textgreater{} 0.

Remark 3. --- If \(g: \mathbf{N} \to \mathbf{N}'\) is a homomorphism of left A-modules, then

\[
\left(1 _ {\mathrm{L} (\mathbf {M})} \otimes g\right) \circ \left(1 _ {\mathrm{L} (\mathbf {M})} \otimes 1 _ {\mathrm{N}}\right) = \left(1 _ {\mathrm{L} (\mathbf {M})} \otimes 1 _ {\mathrm{N}}\right) \circ \left(1 _ {\mathrm{L} (\mathbf {M})} \otimes \mathrm{L} (g)\right),
\]

therefore the diagram

\[
\begin{array}{c c c} \operatorname{Tor} ^ {\mathbf {A}} (\mathbf {M}, \mathbf {N}) & \xrightarrow {\psi_ {\mathbf {M}} (\mathbf {N})} & \mathrm{H} (\mathrm{L} (\mathbf {M}) \otimes_ {\mathbf {A}} \mathbf {N}) \\ \operatorname{Tor} ^ {\mathbf {A}} (1, g) \Big \downarrow & & \Big \downarrow \mathrm{H} (1 \otimes g) \\ \operatorname{Tor} ^ {\mathbf {A}} (\mathbf {M}, \mathbf {N} ^ {\prime}) & \xrightarrow {\psi_ {\mathbf {M}} (\mathbf {N} ^ {\prime})} & \mathrm{H} (\mathrm{L} (\mathbf {M}) \otimes_ {\mathbf {A}} \mathbf {N} ^ {\prime}) \end{array}
\]

is commutative.

Likewise, if \(f: \mathbf{M} \to \mathbf{M}'\) is a homomorphism of right A-modules, we have a commutative diagram:

\[
\begin{array}{c c} \operatorname{Tor} ^ {\mathbf {A}} (\mathrm{M}, \mathrm{N}) & \xrightarrow {\overline {{\psi}} _ {\mathrm{N}} (\mathrm{M})} \mathrm{H} (\mathrm{M} \otimes_ {\mathrm{A}} \mathrm{L} (\mathrm{N})) \\ \operatorname{Tor} ^ {\mathbf {A}} (f, 1) \Big \downarrow & \Big \downarrow \mathrm{H} (f \otimes 1) \\ \operatorname{Tor} ^ {\mathbf {A}} (\mathrm{M} ^ {\prime}, \mathrm{N}) & \xrightarrow {\overline {{\psi}} _ {\mathrm{N}} (\mathrm{M} ^ {\prime})} \mathrm{H} (\mathrm{M} ^ {\prime} \otimes_ {\mathrm{A}} \mathrm{L} (\mathrm{N})). \end{array}
\]

PROPOSITION 6. --- The map (f, g) ↦ Tor \(^{A}\) (f, g) :

\[
\operatorname{Hom} _ {\mathrm{A}} \left(\mathrm{M}, \mathrm{M} ^ {\prime}\right) \times \operatorname{Hom} _ {\mathrm{A}} \left(\mathrm{N}, \mathrm{N} ^ {\prime}\right)\rightarrow \operatorname{Hom} _ {k} \left(\operatorname{Tor} ^ {\mathrm{A}} (\mathrm{M}, \mathrm{N}), \operatorname{Tor} ^ {\mathrm{A}} \left(\mathrm{M} ^ {\prime}, \mathrm{N} ^ {\prime}\right)\right)
\]

is k-bilinear.

Let \(f \in \operatorname{Hom}_{\mathbf{A}}(\mathbf{M}, \mathbf{M}')\) , \(g_1, g_2 \in \operatorname{Hom}_{\mathbf{A}}(\mathbf{N}, \mathbf{N}')\) , \(\lambda_1, \lambda_2 \in k\) . Then the morphisms

\[
\lambda_ {1} (\mathrm{L} (f) \otimes g _ {1}) + \lambda_ {2} (\mathrm{L} (f) \otimes g _ {2}) \quad \text { and } \quad \mathrm{L} (f) \otimes (\lambda_ {1} g _ {1} + \lambda_ {2} g _ {2})
\]

of L(M) ⊗\_A N in L(M) ⊗\_A N′ coincide; by prop. 5 and remark 3, we therefore have

\[
\operatorname{Tor} ^ {\mathbf {A}} \left(f, \lambda_ {1} g _ {1} + \lambda_ {2} g _ {2}\right) = \lambda_ {1} \operatorname{Tor} ^ {\mathbf {A}} \left(f, g _ {1}\right) + \lambda_ {2} \operatorname{Tor} ^ {\mathbf {A}} \left(f, g _ {2}\right).\tag{10}
\]

We reason similarly for the mapping \(f\mapsto\mathrm{Tor}^{\mathbf{A}}(f,g)\) .

COROLLARY. --- Let \(\lambda\in k\) . If \(\lambda\) annihilates M or N, it annihilates Tor \(^{A}\) (M, N).

Indeed, \(\lambda .1_{\mathrm{Tor}(\mathbf{M},\mathbf{N})} = \mathrm{Tor}(\lambda .1_{\mathbf{M}},1_{\mathbf{N}}) = \mathrm{Tor}(1_{\mathbf{M}},\lambda .1_{\mathbf{N}})\) .

PROPOSITION 7. --- Let I and J be two sets, \((\mathbf{M}_{\alpha})_{\alpha \in I}\) a family of right A-modules, \((\mathbf{N}_{\beta})_{\beta \in J}\) a family of left A-modules. The homomorphism

\[
\bigoplus_ {\alpha \in I, \beta \in J} \operatorname{Tor} ^ {A} \left(M _ {\alpha}, N _ {\beta}\right)\rightarrow \operatorname{Tor} ^ {A} \left(\bigoplus_ {\alpha \in I} M _ {\alpha}, \bigoplus_ {\beta \in J} N _ {\beta}\right)
\]

deduced from the canonical homomorphisms \(M_{\alpha} \rightarrow \oplus M_{\alpha}\) and \(N_{\beta} \rightarrow \oplus N_{\beta}\) is bijective. It suffices to prove that for every right module M (resp. every left module N), the canonical homomorphism

\[
\bigoplus_ {\beta \in J} \operatorname{Tor} ^ {A} (M, N _ {\beta}) \rightarrow \operatorname{Tor} ^ {A} (M, \bigoplus_ {\beta \in J} N _ {\beta})
\]

\((\text{resp. } \bigoplus_{\alpha \in I} \text{Tor}^{\mathbf{A}}(\mathbf{M}_{\alpha}, \mathbf{N}) \to \text{Tor}^{\mathbf{A}} (\bigoplus_{\alpha \in I} \mathbf{M}_{\alpha}, \mathbf{N}))\) is bijective. This follows from what precedes, from Proposition 1 of X, p.~28, and from the canonical isomorphisms:

\[
\begin{array}{l} \underset {\beta} {\oplus} (\mathrm{L(M)} \otimes_ {\mathrm{A}} \mathrm {N}_ {\beta}) \to \mathrm{L(M)} \otimes_ {\mathrm{A}} (\underset {\beta} {\oplus} \mathrm {N}_ {\beta}), \\ \underset {\alpha} {\oplus} (\mathrm {M}_ {\alpha} \otimes_ {\mathrm{A}} \mathrm{L(N)}) \to (\underset {\alpha} {\oplus} \mathrm {M}_ {\alpha}) \otimes_ {\mathrm{A}} \mathrm{L(N)}. \end{array}
\]

An analogous argument gives:

PROPOSITION 8. --- Let I (resp. J) be a preordered set filtered to the right, \(((M_{\alpha}), (u_{\alpha' \alpha}))\) (resp. \(((N_{\beta}), (v_{\beta' \beta}))\) ) an inductive system of right (resp. left) A-modules relative to I (resp. J). The homomorphism of graded k-modules

\[
\lim _ {(\alpha , \beta) \in I \times J} \operatorname{Tor} ^ {A} \left(M _ {\alpha}, N _ {\beta}\right)\rightarrow \operatorname{Tor} ^ {A} \left(\lim _ {\alpha \in I} M _ {\alpha}, \lim _ {\beta \in J} N _ {\beta}\right),
\]

deduced from canonical A-homomorphisms \(\mathbf{M}_{\alpha} \to \varinjlim \mathbf{M}_{\alpha}\) and \(\mathbf{N}_{\beta} \to \varinjlim \mathbf{N}_{\beta}\) is bijective.

In particular, taking \(J = I\) and noting that the \((\alpha, \alpha), \alpha \in I\) form a cofinal subset of \(I \times I\) , we obtain:

COROLLARY. --- Let I be a preordered set filtered to the right, \((\mathbf{M}_{i}, u_{ji})\) (resp. \((\mathbf{N}_{i}, v_{ji})\) ) an inductive system of right (resp. left) A-modules relative to I. The homomorphism of graded k-modules

\[
\varinjlim_ {i \in I} \operatorname{Tor} ^ {A} (M _ {i}, N _ {i}) \rightarrow \operatorname{Tor} ^ {A} (\varinjlim_ {i \in I} M _ {i}, \varinjlim_ {i \in I} N _ {i}),
\]

deduced from canonical A-homomorphisms \(\mathbf{M}_i\to \varinjlim \mathbf{M}_i\) and \(\mathbf{N}_j\to \varinjlim \mathbf{N}_j\) is bijective.

Let M be a right A-module, N a left A-module, A° the opposite ring of A, M° the left A°-module underlying M, N° the right A°-module underlying M. We have L(M°) = L(M)° and L(N°) = L(N)°, whence a commutation isomorphism (X, p.~63, prop. 2)

\[
\sigma (\mathrm{L} (\mathrm{M}), \mathrm{L} (\mathrm{N})): \mathrm{L} (\mathrm{M}) \otimes_ {\mathrm{A}} \mathrm{L} (\mathrm{N}) \rightarrow \mathrm{L} (\mathrm{N} ^ {\circ}) \otimes_ {\mathrm{A} ^ {0}} \mathrm{L} (\mathrm{M} ^ {\circ}).
\]

Passing to homology, \(\sigma(\mathrm{L}(\mathbf{M}),\mathrm{L}(\mathbf{N}))\) induces a graded isomorphism of degree 0 \(\sigma_{\mathbf{M},\mathbf{N}}:\operatorname{Tor}^{\mathbf{A}}(\mathbf{M},\mathbf{N})\to\operatorname{Tor}^{\mathbf{A}^{\circ}}(\mathbf{N}^{\circ},\mathbf{M}^{\circ})\) called the commutation isomorphism of torsion products.

Note that \(\sigma_{\mathbf{N}^{\circ},\mathbf{M}^{\circ}}\circ\sigma_{\mathbf{M},\mathbf{N}}=\mathrm{Id}_{\mathrm{Tor}\left(\mathbf{M},\mathbf{N}\right)}\) and that \(\sigma_{\mathbf{M},\mathbf{N}}\) induces on the degree 0 terms the commutation homomorphism of the tensor product. On the other hand, if \(f:\mathbf{M}\to\mathbf{M}'\) and \(g:\mathbf{N}\to\mathbf{N}'\) are homomorphisms of A-modules, we have

\[
\operatorname{Tor} ^ {\mathbf {A} ^ {\circ}} (g, f) \circ \sigma_ {\mathbf {M}, \mathbf {N}} = \sigma_ {\mathbf {M} ^ {\prime}, \mathbf {N} ^ {\prime}} \circ \operatorname{Tor} ^ {\mathbf {A}} (f, g).
\]

